\documentclass[11pt]{article}
\usepackage[margin=1in]{geometry}
\usepackage{amsmath,amssymb}
\usepackage{hyperref}
\usepackage{xcolor}
\title{Dense Exponentials Do Not Force a Uniform Algebra to Be Trivial}
\author{Literature-status note for arXiv:math/0308226}
\date{9 August 2026}

\begin{document}
\maketitle

\begin{center}
\fcolorbox{black}{gray!10}{\parbox{0.9\linewidth}{
\textbf{Status: literature already answered (negative).}
This is an exact later-literature resolution, not a new run proof.}}
\end{center}

\section*{Original conjecture}

In Chapter 5 of \emph{Extensions of Normed Algebras}, PDF page 53,
Dawson asks whether the following condition characterizes trivial uniform
algebras \cite{Dawson}:
\[
\overline{\exp A}=A
\qquad\Longrightarrow\qquad
A=C(X).
\]
Equivalently: must a uniform algebra whose exponentials are dense be trivial?
The thesis reports that no counterexample was then known and separately
highlights the natural metrizable case.

\section*{Decisive later answer}

Izzo restates the exact question as Question 1.3 and says explicitly that the
paper answers it \cite{Izzo}.  Theorem 1.6 asserts that there exists a
nontrivial uniform algebra $A$ for which
\[
\exp A=\{\exp a:a\in A\}
\quad\text{is dense in }A.
\]
The stronger Theorem 1.7 constructs a compact polynomially convex set
$X\subset\mathbb C^2$ of topological dimension one such that $P(X)$ is a
nontrivial strongly regular uniform algebra, is Dirichlet on $X$, and has
dense exponentials.  The paper also proves that $X$ is the maximal ideal
space of $P(X)$.  Therefore the conjectured implication is false.

The supporting authors knew they were answering this problem: the abstract
calls it a 17-year-old question of Dales and Feinstein, and the introduction
labels and answers the dense-exponentials question directly.

\section*{Scope left open}

The counterexample has compact metrizable, simply coconnected,
one-dimensional maximal ideal space.  It is not locally connected.  Thus the
counterexample settles the general conjecture, including the broad
``natural and metrizable'' formulation, but it does not settle stronger
variants imposing local connectedness or requiring maximal ideal space
$[0,1]$.  Izzo explicitly records related planar and locally connected
questions as open.

\section*{Search evidence}

The exact phrase ``dense exponential group'' and the formula
$\overline{\exp A}=A$ were searched in the run indexes, local arXiv corpus,
and bounded web search.  The local source for arXiv:2403.19583 contains the
explicit Question 1.3, Theorems 1.6 and 1.7, and the construction's proof of
nontriviality and dense exponentials.  No originality claim is made here.

\begin{thebibliography}{9}
\bibitem{Dawson}
T. W. Dawson, \emph{Extensions of Normed Algebras}, Ph.D. thesis,
arXiv:math/0308226 (2003), Chapter 5, PDF page 53.

\bibitem{Izzo}
A. J. Izzo, \emph{A nontrivial uniform algebra Dirichlet on its maximal
ideal space}, arXiv:2403.19583, Question 1.3 and Theorems 1.6--1.7.
\end{thebibliography}

\end{document}
